\documentclass[../../../../main.tex]{subfiles}
\begin{document}

\begin{flushright}
\footnotesize\textit{【自動文字起こし・要確認】}
\end{flushright}

\noindent
[\textbf{解}] $A$からの流出速度を $v$，水の単位体積あたりの重さを $w$ とすると，

\begin{enumerate}
    \item[ア] 時刻 $k$ における総水溶液は $vk$
    \item[イ] その重さは
    \[
    \begin{aligned}
    w \int_0^k (1+ae^{-bt}) v \, dt &= \left[ t - \frac{a}{b} e^{-bt} \right]_0^k v w \\
    &= vw \left[ \left( k - \frac{a}{b} e^{-bk} \right) - \left( -\frac{a}{b} \right) \right] \\
    &= vw \left( k + \frac{a}{b} - \frac{a}{b} e^{-bk} \right)
    \end{aligned}
    \]
\end{enumerate}

だから，求める比重は
\[
\frac{vw \left( k + \frac{a}{b}(1 - e^{-bk}) \right)}{vk} \cdot \frac{1}{w} = 1 + \frac{a}{bk} (1 - e^{-bk})
\]

\begin{center}
\small
\begin{tabular}{|p{0.9\textwidth}|}
\hline
\textbf{補足メモ}: \\
「比重\dots{}水に対する相対重さ」という解釈で良いのかな \\
$\Rightarrow \frac{(\text{重さ})}{(\text{総量})}$ が単位あたり重さ，これと水の重さでくらべる \\
\hline
\end{tabular}
\end{center}

\newpage

\end{document}
